\subsection{Passive ordering-call diagnostics}
The completed hybrid records also contain cumulative counters for the
ordering phase. We reconciled these counters across all 27 training runs
and all 35 reused-validation runs, checking each result hash against the
independent audit and the interpreted source code against its collection
protocol. Table~\ref{tab:ordering-counters} reports the retained categories.

\begin{table}[htbp]
\centering
\small
\begin{tabular}{lrrrrr}
\toprule
Study & Runs & Calls & Exit 0 & Timeouts & Other \\
\midrule
TRAIN & 27 & 19,778 & 19,443 & 327 & 8 \\
Reused validation & 35 & 72,388 & 72,234 & 150 & 4 \\
\bottomrule
\end{tabular}
\caption{Passive counters in completed 60-second hybrid runs. Calls are
invocations of the ordering wrapper. Exit 0 records normal process exit;
timeouts record a reported subprocess timeout; other combines nonzero or
absent exit codes. These categories partition wrapper invocations, rather
than validated improvements or feasible final solutions.}
\label{tab:ordering-counters}
\end{table}


The counters require a distinction between a call to the ordering wrapper
and a launched LKH subprocess. A wrapper call can exhaust its allowance
during matrix preparation, yielding no process exit code. A zero exit code
does not by itself show that an improving tour was returned, validated and
accepted before the deadline. Similarly, the other/absent-exit category
cannot distinguish a process failure from a preparation-only invocation.
The retained aggregate records do not supply per-call output-validity or
improvement counts. Final route feasibility and incumbent eligibility are
checked separately.

In training, rd400 accounted for 277 of the 327 reported subprocess timeouts.
Its 735 wrapper invocations therefore had a timeout fraction of 37.7\%.
The corresponding fraction across reused validation was 150 of 72,388
invocations, or 0.21\%. Each of the four tied validation geometries had at
most four such timeouts. Widespread timeouts thus do not describe the
execution status of those tied cases. This observation does not explain
their objective values or establish optimality of their visited supports.
The counters can include a final call returning after the outer deadline;
they are not counts restricted to eligible improvement events. No operator,
budget allocation or screening threshold is changed by this diagnostic.
